\documentclass[12pt]{article}
\usepackage[utf8]{inputenc}
\usepackage[T1]{fontenc}
\usepackage{hyperref}
\title{Fam Language Resources for K-12 Curriculum Development}
\author{Compiled by Claude Sonnet 5.5 (claude-sonnet-5-5) for Andr\'as Kornai}
\date{}
\begin{document}
\maketitle

\section{Introduction}
Fam (ISO 639-3 \texttt{fam}; also called Awol Fam, and by outsiders Konga) is an endangered language of Bali LGA, Taraba State, Nigeria. Sources give fewer than 500 speakers, based on an old (1984) record. Its classification is uncertain; Blench has proposed a divergent Mambiloid affiliation. Very little is documented, and no school use is known.

\section{Primary Educational Resources}
No primers, textbooks or school materials in or for Fam were found in web searches.

\section{Linguistic and Reference Materials}
\begin{itemize}
\item Wikipedia, ``Fam language'': overview of names, location, speaker numbers and classification. \url{https://en.wikipedia.org/wiki/Fam_language}
\item Endangered Languages Project entry for Fam (vitality rating and basic data). \url{https://endangeredlanguages.com/lang/5994}
\end{itemize}

\section{Religious and Community Resources}
\begin{itemize}
\item Joshua Project profile of the Fam people in Nigeria (mission-oriented demographic profile). \url{https://joshuaproject.net/people_groups/11186/NI}
\end{itemize}
No Scripture translation or community organisation was found.

\section{Digital and Audio Resources}
No audio, video, corpus or software resources were found.

\section{Curriculum Development Context}
Fam appears to have no written standard, no teaching materials and a very small, aging speaker base. Any curriculum would need to begin with fresh documentation work with speakers, including an orthography and basic word lists, and the current speaker situation should be verified in the field before planning.
\end{document}